%%=================================================================
%% DISCUSSION: Model Comparison, Neutrino Connection, Limitations
%%=================================================================

\section{Comparison with Alternative Dark Matter Models}
\label{disc:comparison}

Table~\ref{tab:model_comparison} summarizes how our framework compares
with five leading alternatives across the key phenomenological tests.

\begin{table}[h!]
\centering
\small
\caption{Model comparison across key observables. \checkmark = naturally
         explained, $\sim$ = requires tuning, \textsf{X} = fails or
         does not address, --- = not applicable. Only the present model
         simultaneously addresses SIDM small-scale problems, relic density,
         direct detection evasion, and dark energy within a single Lagrangian.}
\label{tab:model_comparison}
\begin{tabular}{l@{\hspace{4pt}}c@{\hspace{4pt}}c@{\hspace{4pt}}c@{\hspace{4pt}}c@{\hspace{4pt}}c@{\hspace{4pt}}c}
\toprule
Observable & CDM & WIMP & FDM & aSIDM & vdSIDM & \textbf{Ours} \\
\midrule
Core-cusp               & \textsf{X} & \textsf{X} & $\sim$ & \checkmark & \checkmark & \checkmark \\
Missing satellites       & $\sim$     & $\sim$     & \checkmark & \checkmark & \checkmark & \checkmark \\
Too-big-to-fail         & \textsf{X} & \textsf{X} & $\sim$ & \checkmark & \checkmark & \checkmark \\
Halo diversity           & \textsf{X} & \textsf{X} & \textsf{X} & $\sim$ & \checkmark & \checkmark \\
Cluster mergers          & \checkmark & \checkmark & \checkmark & $\sim$ & \checkmark & \checkmark \\
Fornax GC survival       & \textsf{X} & \textsf{X} & $\sim$ & $\sim$ & \checkmark & \checkmark \\
Correct relic $\Omega h^2$ & postulated & \checkmark & postulated & $\sim$ & $\sim$ & \checkmark \\
Direct detection safe    & ---        & $\sim$     & ---    & $\sim$ & $\sim$ & \checkmark \\
$\Delta N_{\rm eff}$ ($0.16$--$0.26$) & \checkmark & \checkmark & \checkmark & $\sim$ & $\sim$ & \checkmark \\
Dark energy              & separate   & separate   & separate & separate & separate & \checkmark \\
$H_0$ value             & input      & input      & input  & input  & input  & \textbf{67.4} \\
\# free parameters       & 6          & $\geq 4$   & 1      & 3      & 3      & \textbf{3} \\
\bottomrule
\end{tabular}
\end{table}

We briefly discuss each alternative:

\textbf{$\Lambda$CDM.} Remarkably successful on large scales ($> 1$~Mpc) but
fails on small scales — core-cusp, too-big-to-fail, and diversity problems
persist even with baryonic feedback. The cosmological constant is an input,
not a prediction. Our model reduces to $\Lambda$CDM phenomenology on large
scales (velocity-suppressed $\sigmaT$) while resolving the small-scale
tensions.

\textbf{WIMPs (neutralino, Higgsino, etc.).} Provide the correct relic
density via the ``WIMP miracle,'' but have no self-interaction mechanism
and thus inherit all CDM small-scale problems. Furthermore, they face
growing tension with LZ/XENONnT direct detection null results (except for
finely tuned blind spots). Our model achieves the relic density without
any direct detection signal by construction.

\textbf{Fuzzy dark matter (FDM, $m \sim 10^{-22}$~eV).} Creates cores via
quantum pressure (de~Broglie wavelength $\sim$ kpc). However, Lyman-$\alpha$
forest observations constrain $m > 2\times10^{-20}$~eV~\cite{Irsic2017},
in strong tension with the core sizes needed for dwarf galaxies. FDM also
cannot explain halo diversity (all halos get the same core for given mass)
and does not address dark energy.

\textbf{Constant SIDM (aSIDM, $\sigma/m = \text{const}$).} Solves dwarf-scale
problems but is excluded by cluster mergers if $\sigma/m \gtrsim 1$~cm$^2$/g.
The velocity dependence in our model is automatic — it emerges from the
Yukawa potential, not from an ad hoc form factor.

\textbf{Velocity-dependent SIDM (vdSIDM), generic.}
This is the closest competitor. Many groups have explored Yukawa-mediated SIDM,
typically with Dirac fermions~\cite{Kaplinghat2016,TYZ2013}. Our model differs
in three respects: (i) the fermion is Majorana, producing a distinct
spin-statistics fingerprint (§\ref{ch1:falsifiable});
(ii) the dark sector is fully secluded, eliminating direct detection and CMB
constraints; (iii) the dark energy sector emerges from CP violation in the
same Lagrangian rather than being a separate model.


\section{The \texorpdfstring{$A_4$}{A4} Symmetry and the Neutrino Question}
\label{disc:neutrino}

The $A_4$ discrete symmetry plays a central role in our framework: it governs
the Coleman-Weinberg vacuum structure (§\ref{ch2:cw}) and the CP angle
$\theta_{relic}$.

The same $A_4$ group has been extensively studied in neutrino physics as the
origin of tribimaximal mixing~\cite{Ma2001,AltarelliFeruglio2005}. In those
models, $A_4$ acts on three lepton generations and predicts specific neutrino
mixing angles. In our framework, $A_4$ acts on the \emph{dark sector} — the
dark fermion $\chi$ and mediator $\phi$. The coincidence is suggestive but,
at present, coincidental.

We emphasize what this model \textbf{does not} contain:
\begin{itemize}
\item No neutrino mass mechanism. The secluded dark sector has zero tree-level
      coupling to the Standard Model, so no see-saw or radiative neutrino
      mass is generated.
\item No lepton number violation visible to SM. The Majorana mass of $\chi$
      is in the dark sector; $\chi$ carries no SM quantum numbers.
\item No $\theta_{13}$ correction to dark CP. The neutrino mixing angle
      $\theta_{13}$ does not feed back into the dark sector at tree level.
\end{itemize}

However, the fact that $A_4$ governs \emph{both} neutrino mixing (in the SM
sector) and dark CP violation (in our dark sector) raises the possibility of
a unified $A_4$ UV completion where a single symmetry breaking simultaneously
sets the neutrino mass matrix and the dark sector CW vacuum. This is left
as an explicit open question for future work. The key observable would be a
correlation between the neutrino CP phase $\delta_{CP}$ and the dark sector
relic angle $\theta_{relic} = 18.43^\circ$ — testable with DUNE and Hyper-K.


\section{Relation to Dark QCD and Confinement}
\label{disc:darkqcd}

The dark energy prediction (Chapter~2, §\ref{ch2:h0}) relies on a dark QCD
sector that confines at $\Lambda_d \sim 2.05$~meV. The dark axion $\sigma$ is
the pseudo-Goldstone boson of the spontaneously broken dark chiral symmetry,
with mass $m_\sigma = \Lambda_d^2/f \sim H_0$ via the Gell-Mann--Oakes--Renner
relation.

This sector is not added by hand — it is the non-perturbative completion of
the $A_4$-symmetric dark Yukawa interaction. The QCD-like dynamics provides:
(i) confinement of dark color (eliminating free dark quarks as a relic);
(ii) the instanton-generated potential for $\sigma$;
(iii) the mass gap that protects $\sigma$ from radiative corrections.

The coincidence $\Lambda_d \sim {\rm meV} \sim H_0$ appears fine-tuned but is
in fact a consequence of the Friedmann equation: $V_\sigma = f^2 m_\sigma^2/2 = \rho_\Lambda$
determines $m_\sigma$ once $f \sim 0.2\,M_{\rm Pl}$ is set by Planck-scale dynamics.
It is no more fine-tuned than the electroweak scale in the Standard Model — both
are set by symmetry breaking scales.


\section{Limitations and Open Problems}
\label{disc:limitations}

We list the principal limitations of the present work and their status:

\begin{enumerate}
\item \textbf{Mediator overclosure (resolved).}
      The $\phi$ relic abundance is depleted by cannibal processes
      $3\phi \leftrightarrow 2\phi$, enabled by the cubic self-coupling $\mu_3$
      in the scalar potential (§\ref{ch3:cannibal}). A dedicated Boltzmann
      calculation shows that $\mu_3/m_\phi > 0.85$ suffices to bring
      $\Omega_\phi < \Omega_{\rm CDM}$, with the full perturbative window
      $0.85 < \mu_3/m_\phi < 3.54$. The dark sector temperature at BBN
      remains $T_{\rm SM} > 1$~MeV, preserving nucleosynthesis.
      This resolves the overclosure problem identified in Chapter~1.

\item \textbf{Mixed Majorana coupling.}
      The Chapter~3 decomposition assumes a pure scalar coupling $y_s$ for
      SIDM at $\theta = 0$ and pure pseudoscalar $y_p$ at $\theta = \pi/2$.
      In reality, the relic angle $\theta_{relic} = 18.43^\circ$ means the
      physical coupling is a mixture:
      $y_{\rm phys} = y\cos\theta_{relic} + iy\sin\theta_{relic}\gamma^5$.
      The interference between scalar and pseudoscalar exchange modifies
      the self-interaction cross section at the $\sim 10\%$ level. A full
      calculation of $\sigmaT$ with the mixed $(y_s + iy_p\gamma^5)$ vertex
      is needed.

\item \textbf{Spiral galaxy fits.}
      The RAR analysis uses a simple SIDM core model without adiabatic
      contraction (Gnedin et al.~2004). For spiral galaxies like NGC~2403
      and NGC~3198, the baryonic gravitational potential modifies the DM
      density profile significantly. A self-consistent calculation with
      adiabatic contraction could shift $\sigmaT/m$ by up to $\sim 30\%$.

\item \textbf{DESI DR2 equation of state (largely resolved).}
      The analytic slow-roll estimate gives $w_0 \approx -0.727$, but the
      full numerical Friedmann + Klein-Gordon integration yields
      $w_0 = -0.981$ (§\ref{ch3:bestfit}), since the field remains frozen
      near $\theta_i \approx \pi/2$ at the present epoch. This places the
      model at $2.4\sigma$ from the DESI central value ($w_0 = -0.827$)
      \emph{on the $\Lambda$CDM side}, and indistinguishable from $\Lambda$CDM
      given current SN~Ia (Pantheon+, $\Delta\chi^2 = -9.1$) and BAO data.
      The predicted departure $1 + w_0 \approx 0.02$ is a sharp target for
      next-generation surveys (DESI DR4, Euclid, Roman).

\item \textbf{MCMC convergence at high $\lambda$.}
      The MCMC sampler explores $\lambda$ up to 49 (MAP point), where
      the VPM solver requires $\ell_{\max} \sim 80$ partial waves. The
      convergence diagnostics ($\hat{R} < 1.01$, $\tau_{\max} = 75$) are
      satisfactory, but a production run with $> 3000$ steps and additional
      QA (one-sided $\chi^2$ treatment, $\lambda > 50$ cutoff) is desirable.

\item \textbf{No UV completion.}
      The model is an effective field theory. The Yukawa coupling, $A_4$
      symmetry, and dark QCD confinement are postulated at the Lagrangian level.
      A UV completion (e.g., embedding in a string compactification or a
      composite Higgs-like framework) is not provided.
      We note, however, that the dark confinement scale
      $\Lambda_d \sim 2$~meV arises naturally if the dark sector
      couples to a hidden sector that breaks at
      $\Lambda_{\rm UV} \sim \mathcal{O}(\text{TeV})$ through
      Planck-suppressed operators:
      $\Lambda_d \sim \Lambda_{\rm UV}^2 / M_{\rm Pl} \sim
      (3\;\text{TeV})^2 / (2.4\times10^{18}\;\text{GeV})
      \sim 4\;\text{meV}$,
      in the correct range. This suggests that the meV scale
      is not accidental but may reflect gravity-mediated
      transmission of TeV-scale breaking to the dark sector ---
      a direction left for future investigation.
\end{enumerate}

None of these limitations invalidate the existing results; all cross sections
and relic density calculations are exact within the stated model. The
limitations define the boundary of the present analysis and the agenda for
subsequent work.


\section{Conclusions}
\label{disc:conclusions}

We have presented a unified framework for the dark sector built from a single
Lagrangian with three free parameters $(m_\chi, m_\phi, \alpha)$:

\begin{enumerate}
\item A Majorana fermion $\chi$ interacting via a secluded scalar mediator $\phi$
      provides velocity-dependent self-interactions that resolve the core-cusp,
      too-big-to-fail, and halo diversity problems (Chapter~1).

\item The same Lagrangian, evaluated as a path integral, proves that the VPM
      scattering ODE is the exact fluctuation determinant (no approximation),
      that the CW potential selects a CP-violating vacuum, and that dark QCD
      confinement yields $H_0 = 67.4$~km/s/Mpc, consistent with Planck (Chapter~2).

\item The CP phase, promoted to a dynamical dark axion $\sigma$, provides an
      ``electromagnetic duality'' of the dark sector: the T-even component
      $y_s$ drives SIDM clustering while the T-odd component $y_p$ sources
      the observed vacuum energy (Chapter~3). Numerical integration of the
      Friedmann + Klein-Gordon equations yields $H_0 = 67.39$~km/s/Mpc,
      $\Omega_{\rm DE} = 0.686$, $w_0 = -0.981$, and $w_a = -0.028$ ---
      indistinguishable from $\Lambda$CDM with current data.
      Confrontation with Pantheon+ SN~Ia ($\Delta\chi^2 = -9.1$)
      and BAO ($\chi^2 = 5.4/9$) confirms competitive or better fit
      relative to the concordance model.
\end{enumerate}

The model is tested against 13 astrophysical observables spanning $v = 12$--$4700$~km/s,
achieving $\chi^2_{\min}/\nu = 0.12$ (unconstrained) and $\chi^2_{\rm relic}/\nu = 0.38$
(relic-constrained). All 17 relic-viable benchmark points satisfy every observational
constraint simultaneously. The numerical pipeline is validated by a six-part cross-check
suite including Born-limit agreement, unitarity, transfer$\equiv$elastic identity,
literature comparison, and blind stress tests. The mediator overclosure problem ---
previously a fatal limitation --- is resolved via the cannibal mechanism
$3\phi \to 2\phi$ with tree-level coupling $\mu_3/m_\phi > 0.85$.

The framework makes falsifiable predictions: a Majorana spin-statistics fingerprint
at $v \sim 40$--$95$~km/s (observable in intermediate dSphs), no SMBH seed formation
($t_{\rm gc}/t_{\rm univ} > 10^4$), $w_0 \approx -0.98$ (testable with DESI DR4/Euclid),
and a universal relic angle $\theta_{\rm relic} = 18.43^\circ$ independent of
$(m_\chi, m_\phi, \alpha)$.

All code, data, and cross-check scripts are publicly available~\cite{SIDMcode2026}.
